

\documentclass[letter, 11pt]{article}

    %%%%%%%%%%%%%%%%%%%%%%%%%%%%%%
    \oddsidemargin  0.0in
    \evensidemargin 0.0in
    \textwidth      6.8in
    \headheight     28pt
    \topmargin      0.0in
    \textheight=9.0in
    \parindent=0in
    %%%%%%%%%%%%%%%%%%%%%%%%%%%%%%

\renewcommand{\familydefault}{\sfdefault} % Police 

\usepackage[dvipsnames]{xcolor}
\usepackage{amsthm,amsfonts,amssymb,amsmath, enumitem}
\usepackage{graphicx}   %To include graphics
\usepackage[small,compact]{titlesec}

\usepackage{textcomp}
\newcommand{\mytilde}{$\sim$}
\usepackage{fancyvrb}
\fvset{framesep=1mm,fontfamily=courier,framerule=.3mm,numbersep=1mm,commandchars=\\\{\}}

\usepackage[T1]{fontenc}\usepackage{lmodern}\usepackage[utf8]{inputenc}

\usepackage[french]{babel}

\usepackage{fancyhdr}
\pagestyle{fancy}

\usepackage{version} % Pour faire un texte troué pour les étudiants
%\includeversion{etudiant}\excludeversion{professeur} %version etudiant
\includeversion{professeur}\excludeversion{etudiants} %version prof
\newcommand{\etudiant}[1]{\begin{etudiants}#1 \end{etudiants}}
\newcommand{\prof}[1]{\begin{professeur}\color{blue} {\bf Solution :} #1 \end{professeur}}


\fancyhead[L]{STT-2100 \\ Régression}
\fancyhead[R]{A.-S. Charest, E. Reny-Nolin, A.-A. Vallée \\ Université Laval}


\begin{document}

\vspace*{0.01 cm}
\begin{center}
{\Large Inférence  et prévision en RLS avec R }
\end{center}

Les mots en rouge indiquent les noms choisis par l'utilisateur.\\

Pour créer un objet de type "modèle linéaire" :
\begin{Verbatim}[commandchars=\\\{\}]
{\textcolor{red}{mod}}<- lm({\textcolor{red}{varY}} \mytilde {\textcolor{red}{varX}}, data = {\textcolor{red}{nomDuJeuDeDonnees}})
\end{Verbatim}

\subsection*{Intervalle de confiance de niveau $1-\alpha=0,90$ pour les paramètres:}
\begin{Verbatim}
confint({\textcolor{red}{mod}}, level = {\textcolor{red}{0.90}})
\end{Verbatim}

\subsection*{Test d’hypothèses bilatéral avec $c=0$, soit $H_0:\beta_j=0$ contre $H_0:\beta_j\neq 0$  :}
\begin{Verbatim}
summary({\textcolor{red}{mod}}) {\textcolor{ForestGreen}{#Vous avez les statistiques observées et les valeurs p}}
\end{Verbatim}

\subsection*{Tableau d’analyse de variance :}
\begin{Verbatim}
anova({\textcolor{red}{mod}}) {\textcolor{ForestGreen}{#Vous avez le coefficient de détermination et l'estimation de }}
{\textcolor{ForestGreen}{#la variance des erreurs}}
\end{Verbatim}

\subsection*{Prédiction pour toutes les observations dans le jeu de données :}
\begin{Verbatim}
predict({\textcolor{red}{mod}})
\end{Verbatim}

\subsection*{Prédiction pour d’autres observations :}

Créer d'abord un nouveau jeu de données contenant les valeurs de $x$ pour lesquelles on souhaite une prédiction, disons 12,  20 et 54. Attention : ce nouveau jeu doit être un data frame avec les mêmes noms que celui utilisé dans \texttt{lm()}, et aucune valeur spécifiée pour $Y$.
\begin{Verbatim}[commandchars=\\\{\}]
{\textcolor{red}{nouveau}} <- data.frame( {\textcolor{red}{varX}} = c(12, 20, 54))
\end{Verbatim}
On peut ensuite calculer des prédictions basées sur le modèle: 
\begin{Verbatim}
predict({\textcolor{red}{mod}}, new = {\textcolor{red}{nouveau}})
\end{Verbatim}


\subsection*{Intervalle de confiance/prévision pour les prédictions :
}
\begin{Verbatim}
predict({\textcolor{red}{mod}}, new = {\textcolor{red}{nouveau}}, interval = "confidence") {\textcolor{ForestGreen}{#Pour la moyenne}}
predict({\textcolor{red}{mod}}, new = {\textcolor{red}{nouveau}}, interval = "prediction") {\textcolor{ForestGreen}{#Pour une prévision}}
{\textcolor{ForestGreen}{#Voir help(predict.lm) pour plus de détails.}}
\end{Verbatim}



\end{document}

